\documentclass[sigconf]{acmart}

\settopmatter{printacmref=false}
\renewcommand\footnotetextcopyrightpermission[1]{}
\setcopyright{none}

\acmConference[SBES 2026]{40th Brazilian Symposium on Software Engineering}{September 8--11, 2026}{São Paulo, SP, Brazil}

\renewcommand\abstractname{Resumo}
\renewcommand\keywordsname{Palavras-chave}
\renewcommand\refname{Referências}

\AtBeginDocument{
  \providecommand\BibTeX{{Bib\TeX}}
}

\usepackage{booktabs}
\usepackage{listings}
\usepackage{xcolor}
\usepackage{url}
\usepackage{mdframed}

\newmdenv[
  leftline=true,
  rightline=false,
  topline=false,
  bottomline=false,
  linewidth=3pt,
  linecolor=blue!50!black,
  backgroundcolor=blue!3!white,
  innerleftmargin=8pt,
  innerrightmargin=8pt,
  innertopmargin=4pt,
  innerbottommargin=4pt,
  skipabove=6pt,
  skipbelow=6pt
]{promptbox}
\usepackage{iftex}
\iftutex
  \usepackage{fontspec}
\else
  \usepackage[utf8]{inputenc}
  \usepackage[T1]{fontenc}
\fi
\usepackage[brazil]{babel}
\usepackage{hyperref}
\usepackage{tikz}
\usetikzlibrary{calc}

\begin{document}

\title{Dívida Técnica Arquitetural em Sistemas Gerados por IA Generativa: um estudo experimental sobre MVPs e escalabilidade}

\author{Anne Esther Lins Figueirôa}
\affiliation{
  \institution{Instituto de Tecnologia e Liderança -- INTELI}
  \city{São Paulo}
  \country{Brazil}
}
\email{anne.figueiroa@inteli.edu.br}

\author{Reginaldo Arakaki}
\affiliation{
  \institution{Instituto de Tecnologia e Liderança -- INTELI}
  \city{São Paulo}
  \country{Brazil}
}
\email{reginaldo.arakaki@inteli.edu.br}

\author{Fabiana Martins de Oliveira}
\affiliation{
  \institution{Universidade de São Paulo -- USP}
  \city{São Paulo}
  \country{Brazil}
}
\email{fabiana.martins@usp.br}

\renewcommand{\shortauthors}{Figueirôa et al.}
\renewcommand{\shorttitle}{Dívida Técnica Arquitetural em Soluções Geradas por IA}

\renewcommand\abstractname{Abstract}
\begin{abstract}
Generative AI (GenAI) tools have accelerated MVP development in startups, enabling
founders without technical backgrounds to build functional systems via natural language
prompts. However, little is known about the architectural quality of such systems under
real-world scale conditions. This paper presents a controlled experimental study
evaluating the architectural technical debt (ATD) introduced by GenAI tools compared
to architectures optimized via explicit architectural prompts. Based on a persona
derived from semi-structured interviews with real startup founders, a backend was
implemented in two scenarios on the same stack (Node.js/PostgreSQL): Scenario A,
generated by Lovable with business prompts lacking architectural specification; and
Scenario B, refactored with four explicit architectural decisions. Results showed that
Scenario A completed the protocol without HTTP errors but with severe latency
degradation (mean 1,880\,ms, p95 threshold violated), while Scenario B responded
in 30\,ms on average, meeting all criteria: a 98.4\% difference associated with the
architectural decisions applied to an identical technology stack, data, and
infrastructure. The generation tool (a no-code builder in Scenario A vs.\ an
LLM-based coding assistant in Scenario B) and construction mode (autonomous
generation vs.\ guided refactoring) differ between scenarios and are addressed
as limitations.
As a practical contribution, a four-component architectural prompt model is proposed,
enabling founders to specify scalability constraints that GenAI models can translate
into correct implementation patterns.
\end{abstract}

\keywords{Dívida Técnica Arquitetural, Inteligência Artificial Generativa, Escalabilidade,
  Prompt Engineering, Startup, Teste de Carga, K6, Decisão Arquitetural, Táticas Arquiteturais}

\frenchspacing
\maketitle

\section*{Resumo}
Ferramentas de Inteligência Artificial Generativa (IAG) aceleraram o desenvolvimento
de Produtos Mínimos Viáveis (MVPs) em startups, permitindo que fundadores sem formação
técnica construam sistemas funcionais via prompts em linguagem natural. Contudo, pouco
se sabe sobre a qualidade arquitetural desses sistemas quando submetidos a condições
reais de escala. Este artigo apresenta um estudo experimental que avalia a dívida técnica
arquitetural (DTA) introduzida por ferramentas de IAG em comparação com arquiteturas
otimizadas via prompts arquiteturais explícitos. A partir de uma persona baseada em
entrevista semiestruturada com fundadores de uma startup real, implementou-se um servidor
de aplicação em dois cenários sobre a mesma pilha tecnológica (Node.js/PostgreSQL):
Cenário~A, gerado pelo Lovable com prompts de negócio sem especificação arquitetural;
e Cenário~B, obtido pela refatoração do código do Cenário~A via prompts arquiteturais
explícitos enviados a um assistente de codificação baseado em LLM (Claude). A ferramenta
de geração e o modo de construção diferem entre os cenários, o que é discutido como
limitação do estudo. A análise estática
revelou que o Cenário~A executa 109 consultas ao banco por requisição no dashboard do
gestor (padrão N+1), contra uma única consulta no Cenário~B (JOIN). Testes de carga
com 100 usuários virtuais simultâneos por 3 minutos evidenciaram latência média de
1.880\,ms no Cenário~A contra 30\,ms no Cenário~B (−98,4\%), com 0\% de erro em
ambos. Como contribuição prática, propõe-se um modelo de prompt arquitetural de quatro
componentes que permite a fundadores especificar restrições de escalabilidade que
modelos de IAG consigam traduzir em padrões de implementação corretos.


%% ─────────────────────────────────────────────────────────────────────────────
\section{Introdução}
%% ─────────────────────────────────────────────────────────────────────────────

Em 2024, 60\% das startups utilizavam ferramentas \textit{no-code}/\textit{low-code}
e 74\% dos empreendedores incorporavam IAG como componente central de seus
produtos~\cite{RudysAI2024,HubSpot2024}. O número de startups fundadas por um único
empreendedor cresceu 23\% no mesmo período, viabilizadas por ferramentas como Lovable,
v0 e Replit Agent, que permitem construir MVPs em horas via descrições em linguagem
natural, consolidando o \textit{vibe coding}: desenvolvimento guiado inteiramente
por prompts, sem especificação de requisitos técnicos~\cite{HubSpot2024}.\footnote{\textit{Vibe coding}: termo cunhado por Andrej Karpathy em 2025 para descrever o desenvolvimento guiado inteiramente por descrições em linguagem natural, sem escrita manual de código.}

Embora eficazes na validação funcional, essas ferramentas abstraem decisões
arquitetônicas críticas: estrutura de acesso a dados, gerenciamento de conexões,
estratégias de cache e processamento assíncrono. A história confirma os riscos:
o Twitter colapsou ao tentar escalar seu monólito em Ruby on Rails; o Uber documentou
migração dolorosa para microsserviços~\cite{Newman2021}.

Pesquisas recentes focam a correção sintática de código gerado por
IAG~\cite{Perry2023,Liu2023,Shypula2023}, mas há escassez de estudos sobre
\textit{integridade arquitetural sistêmica}, a qualidade das decisões de design
que emergem do processo de geração por Inteligência Artificial Generativa (IAG). Esta lacuna é especialmente crítica no
contexto de startups: diferentemente de times de engenharia com práticas estabelecidas
de revisão de código e testes de desempenho, fundadores que adotam IAG como substituto
de equipe técnica raramente dispõem dos instrumentos para detectar Dívida Técnica Arquitetural (DTA) antes que ela
se manifeste em produção, como colapso sob pico de usuários, degradação de
latência com o crescimento da base de dados, ou timeout em funcionalidades
transversais. O presente trabalho aborda diretamente essa lacuna.

\textbf{Hipótese central:} a qualidade arquitetural do código gerado por Inteligência
Artificial Generativa está fortemente associada ao nível de especificação arquitetural
presente nos prompts. Quanto maior a explicitação das restrições arquiteturais, maior
a probabilidade de geração de soluções compatíveis com requisitos de escalabilidade.

\textbf{Questão de Pesquisa:} Como o nível de especificação arquitetural presente em
prompts influencia a dívida técnica arquitetural e a escalabilidade de sistemas
gerados por IAG?

\textbf{Objetivo geral:} avaliar experimentalmente essa hipótese via dois cenários
sobre a mesma pilha tecnológica Node.js/PostgreSQL: \textbf{Cenário~A}, gerado por IAG sem
especificação arquitetural, e \textbf{Cenário~B}, refatorado com quatro decisões
explícitas. Os objetivos específicos são: (I)~mapear táticas de
escalabilidade~\cite{Bass2021} e refatoração~\cite{Fowler2018};
(II)~implementar o Cenário~A via Lovable; (III)~implementar o Cenário~B com
JOIN, pool otimizado, \textit{cache} e fila assíncrona; (IV)~analisar estaticamente
o código-fonte e submeter ambos a testes de carga~\cite{grafana2024}; e
(V)~comparar os padrões arquiteturais e seu impacto na escalabilidade via
ATAM~\cite{Kazman2000}.

%% ─────────────────────────────────────────────────────────────────────────────
\section{Fundamentação Teórica}
%% ─────────────────────────────────────────────────────────────────────────────

\subsection{Dívida Técnica e Sua Dimensão Arquitetural}

O conceito de Dívida Técnica foi cunhado por Cunningham~\cite{Cunningham1992} e
sistematizado por Fowler~\cite{Fowler2018} no ``Quadrante da Dívida Técnica'',
classificando-a em prudente/imprudente e deliberada/inadvertida. Em startups que
adotam IAG predomina a dívida \textit{impru\-dente-ina\-dver\-tida}: a equipe não percebe
o passivo arquitetural por falta de repertório técnico para identificá-lo.
Fernández-Sánchez et al.~\cite{FernandezSanchez2020} identificam que a DTA é a de
maior impacto no ciclo de vida dos sistemas: ao contrário da dívida de código
(localizada), a DTA é sistêmica: uma decisão equivocada sobre padrão de acesso
a dados contamina toda a base de código. A norma ISO/IEC 25010~\cite{ISO25010}
fornece o referencial normativo para mensuração objetiva da DTA.

\subsection{Escalabilidade Como Atributo Arquitetural}

Bass, Clements e Kazman~\cite{Bass2021} distinguem escalabilidade horizontal
(adição de instâncias) e vertical (aumento de recursos por instância). Sistemas
gerados sem especificação arquitetural tendem a incorporar estado compartilhado em
memória e chamadas síncronas em cascata, que inviabilizam o escalonamento
horizontal. Kleppmann~\cite{Kleppmann2017} descreve os compromissos (\textit{tradeoffs}) centrais (teorema
CAP, latência vs.\ \textit{throughput}) e confirma que as estratégias de menor complexidade
imediata: consultas síncronas, sem cache, processamento sequencial, que correspondem
ao maior custo sob escala. O ATAM (\textit{Architecture Tradeoff Analysis
Method})~\cite{Kazman2000}, desenvolvido pelo SEI da Carnegie Mellon University,
identifica pontos de sensibilidade (decisões que afetam criticamente um único
atributo de qualidade) e compromissos (\textit{tradeoffs}), ou seja, decisões que afetam múltiplos atributos de forma conflitante, sendo adotado neste estudo como instrumento de análise
qualitativa-comparativa entre os cenários.

\subsection{Avaliação RM-ODP dos Cenários}
\label{sec:rmodp}

Ferramentas como o Lovable traduzem requisitos de negócio em decisões arquiteturais
autônomas, sem que o fundador perceba as implicações. O arcabouço RM-ODP (\textit{Reference Model for Open Distributed Processing}) distingue cinco
\textit{viewpoints} para especificação de sistemas distribuídos: \textit{enterprise}
(objetivos de negócio), \textit{information} (semântica dos dados),
\textit{computational} (decomposição funcional), \textit{engineering} (mecanismos
de distribuição) e \textit{technology} (implementação). Fundadores operam
exclusivamente no viewpoint \textit{enterprise}; a IAG realiza a tradução para os
demais de forma autônoma, com qualidade arquitetural variável. A
Figura~\ref{fig:rmodp} compara os dois cenários nos cinco viewpoints, usando escala
ordinal 1--5 adaptada da ISO/IEC 25010~\cite{ISO25010}: 1 indica ausência de
tratamento do atributo; 5 indica cobertura completa. A avaliação tem caráter
qualitativo-comparativo, realizada por um dos autores via inspeção estática do
código e análise dos artefatos de teste.

A Figura~\ref{fig:rmodp} evidencia que as maiores diferenças entre os cenários
concentram-se nos viewpoints \textit{Computational} e \textit{Engineering}, justamente
aqueles responsáveis pelas decisões arquiteturais relacionadas à escalabilidade,
distribuição e gerenciamento dos recursos do sistema. Essa observação reforça a
hipótese de que a ausência de restrições arquiteturais explícitas nos prompts
compromete principalmente os atributos de qualidade associados à arquitetura de
software, enquanto os viewpoints \textit{Enterprise} e \textit{Information} permanecem
equivalentes entre os cenários, pois ambos partem dos mesmos requisitos de negócio.

% Cenário A: Enterprise=4, Information=4, Computational=2, Engineering=1, Technology=2
% Cenário B: Enterprise=4, Information=4, Computational=4, Engineering=4, Technology=4
\begin{figure}[t]
\centering
\begin{tikzpicture}[scale=0.78]
  \def\maxval{5}
  \def\N{5}
  \def\R{2.2cm}

  \foreach \i/\lbl/\ang in {
    1/Enterprise/90,
    2/Information/18,
    3/Technology/-54,
    4/Engineering/-126,
    5/Computational/162}
  {
    \draw[gray!40, thin] (0,0) -- (\ang:\R);
    \node[font=\tiny, align=center] at (\ang:{\R+0.42cm}) {\lbl};
  }

  \foreach \lvl in {1,2,3,4,5}{
    \pgfmathsetmacro{\r}{\lvl/\maxval*2.2}
    \draw[gray!30, thin]
      (90:{\r cm}) --
      (18:{\r cm}) --
      (-54:{\r cm}) --
      (-126:{\r cm}) --
      (162:{\r cm}) -- cycle;
    \node[font=\tiny, gray, left=1pt] at (90:{\r cm}) {\lvl};
  }

  \newcommand{\radarpoint}[2]{%
    (#1:{#2/\maxval*2.2 cm})%
  }

  \draw[red!70!black, thick, dashed, fill=red!10, fill opacity=0.35]
    \radarpoint{90}{4} --
    \radarpoint{18}{4} --
    \radarpoint{-54}{2} --
    \radarpoint{-126}{1} --
    \radarpoint{162}{2} -- cycle;
  \foreach \ang/\val in {90/4, 18/4, -54/2, -126/1, 162/2}
    \fill[red!70!black] \radarpoint{\ang}{\val} circle (1.8pt);

  \draw[green!50!black, thick, fill=green!15, fill opacity=0.40]
    \radarpoint{90}{4} --
    \radarpoint{18}{4} --
    \radarpoint{-54}{4} --
    \radarpoint{-126}{4} --
    \radarpoint{162}{4} -- cycle;
  \foreach \ang/\val in {90/4, 18/4, -54/4, -126/4, 162/4}
    \fill[green!50!black] \radarpoint{\ang}{\val} circle (1.8pt);

  \draw[red!70!black, thick, dashed]
    (-2.1cm,-2.7cm) -- (-1.4cm,-2.7cm)
    node[right, font=\tiny, text=black] {Cenário A};
  \draw[green!50!black, thick]
    (0.2cm,-2.7cm) -- (0.9cm,-2.7cm)
    node[right, font=\tiny, text=black] {Cenário B};
\end{tikzpicture}
\caption{Avaliação RM-ODP: Cenário~A (\textit{baseline} IAG, linha tracejada) vs.\ Cenário~B
  (arquitetura refatorada, linha sólida) nos cinco viewpoints. Escala 1--5
  adaptada da ISO/IEC 25010; avaliação qualitativa-comparativa (ver texto).}
\label{fig:rmodp}
\Description{Gráfico de radar com cinco eixos (Enterprise, Information, Computational, Engineering, Technology) comparando Cenário A e Cenário B. Cenário A (linha tracejada) pontua 4 em Enterprise e Information, e 1--2 nos demais. Cenário B (linha sólida) pontua 4 em todos os eixos.}
\end{figure}

\subsection{Estado da Arte e Trabalhos Relacionados}
\label{sec:trabalhos_relacionados}

No contexto de startups, a adoção de IAG como substituto de equipe técnica é
documentada por Ries~\cite{Ries2012} e amplificada pelo fenômeno do \textit{vibe
coding}: o ato de descrever um produto em linguagem natural e deixar que a IAG
tome todas as decisões de implementação. Esse modelo de desenvolvimento apresenta
uma assimetria crítica: a IAG é treinada em repositórios de código orientados à
funcionalidade e tutoriais de introdução, que raramente incluem exemplos de
otimização de desempenho sob carga ou de dimensionamento de infraestrutura para
escala~\cite{Nascimento2024}. Como resultado, os sistemas gerados tendem a replicar
os padrões mais frequentes nos dados de treinamento, que são os padrões mais
simples, não os mais escaláveis.

A literatura sobre qualidade de código gerado por IAG concentra-se em dois eixos.
O primeiro aborda \textbf{correção funcional}: Perry et al.~\cite{Perry2023}
demonstraram que a qualidade do código depende criticamente do nível de especificação
no prompt; Liu et al.~\cite{Liu2023} identificaram degradação em problemas com
múltiplos componentes interdependentes; Shypula et al.~\cite{Shypula2023}
introduziram o conceito de ``lacuna de desempenho sistêmico'' para descrever a
diferença entre código funcionalmente correto e código eficiente sob operação real.
O segundo eixo aborda \textbf{requisitos não-funcionais}: Nascimento et
al.~\cite{Nascimento2024} identificaram lacunas sistemáticas na geração desses
requisitos por modelos generativos; Vogelsang~\cite{Vogelsang2024} argumenta que
requisitos precisam ser enriquecidos com decisões de design antes de serem usados
como prompts, denominando esse processo trabalho de Engenharia de Requisitos (RE) fundamental irredutível à
automação. O presente trabalho avança ao avaliar a DTA no nível sistêmico, via
testes de carga em condições controladas, e ao propor um modelo operacional de
prompt arquitetural validado empiricamente.

%% ─────────────────────────────────────────────────────────────────────────────
\section{Metodologia}
%% ─────────────────────────────────────────────────────────────────────────────

Esta Iniciação Científica foi desenvolvida individualmente pela primeira autora no
INTELI (orientador: Prof.\ Arakaki; coorientadora: Prof.\textsuperscript{a} Fabiana
Oliveira/USP), adotando Pesquisa Experimental Aplicada em ambiente laboratorial
controlado (\textit{experimentação em Software})~\cite{Wohlin2012}. Artefatos
disponíveis em repositório
público\footnote{\url{https://github.com/anneestherlf/genai-architectural-debt}}~\cite{GenAIDebt2026}.

\subsection{Coleta de Dados}

A entrevista teve caráter exploratório-formativo: seu objetivo não foi gerar
dados quantitativos, mas ancorar a persona sintética \textit{Onboard} — e, por
extensão, os Cenários A e B — em padrões reais de adoção de IAG observados em
uma startup, garantindo validade ecológica ao desenho experimental controlado
descrito na Seção~3.3.

Foi conduzida uma entrevista semiestruturada (30–45 min) com dois fundadores de startup pré-seed, selecionados por amostragem intencional. Dois critérios guiaram essa seleção: o produto precisava ter sido construído integralmente via IAG, sem envolvimento de engenheiros de software, e a startup precisava estar em estágio pré-seed, momento em que decisões arquiteturais tendem a ser mais implícitas e menos revisadas~\cite{Ries2012}. Optou-se pela abordagem semiestruturada porque ela combina cobertura consistente de temas predefinidos — contexto da empresa, adoção de IAG, decisões arquiteturais tomadas ou omitidas e problemas identificados em produção — com a flexibilidade necessária para explorar problemas emergentes relatados espontaneamente pelos fundadores, o que se mostrou essencial para captar sintomas de DTA que os próprios entrevistados não sabiam nomear tecnicamente.

Os participantes foram pseudonimizados e assinaram Termo de Consentimento Livre e Esclarecido (TCLE); os dados foram analisados por meio da análise temática de Braun e Clarke~\cite{Braun2006}. Os cinco padrões emergentes dessa análise (P1–P5, Tabela~\ref{tab:padroes}) formam a base do restante da metodologia: eles fundamentam a construção da persona Onboard (Seção~3.2), orientam a identificação dos sinais de alerta observáveis por fundadores sem formação técnica (Tabela~\ref{tab:sinais}, Seção~5.1) e, por fim, direcionam as quatro refatorações arquiteturais aplicadas no Cenário~B (Seção~5.3).

\subsection{Construção da Persona}

A análise temática gerou a persona \textit{Onboard}: plataforma SaaS
(\textit{Software as a Service}) multi-tenant de onboarding corporativo, construída
com Lovable e Supabase~\cite{Supabase2024} por dois empreendedores sem formação em
Engenharia de Software. A Tabela~\ref{tab:padroes} apresenta os cinco padrões
emergentes da codificação temática.

\begin{table}[t]
  \caption{Padrões de desenvolvimento da persona Onboard (P1--P5), derivados por
           codificação temática indutiva~\cite{Braun2006}.}
  \label{tab:padroes}
  \small
  \begin{tabular}{p{0.9cm}p{2.3cm}p{2.4cm}}
    \toprule
    \textbf{Padrão} & \textbf{Descrição} & \textbf{Implicação arquitetural} \\
    \midrule
    \textbf{P1} & Prompts sem parâmetros arquiteturais & Decisões de design delegadas à IAG \\
    \textbf{P2} & Adoção guiada por mentores, não por critérios técnicos & Inconsistências entre camadas geradas \\
    \textbf{P3} & Aceitação acrítica dos outputs da IAG & DTA inadvertida acumulada \\
    \textbf{P4} & Arquitetura do protótipo permanece em produção & Escalabilidade nunca testada formalmente \\
    \textbf{P5} & Reconhecimento retrospectivo da lacuna de ES & Consciência tardia do custo acumulado \\
    \bottomrule
  \end{tabular}
\end{table}

\subsection{Cenários Experimentais}

O experimento implementa o módulo de trilhas de aprendizado da Onboard, componente
de maior carga computacional, representativo do padrão Fan-out\footnote{Múltiplos usuários acessando conteúdo personalizado simultaneamente, gerando carga proporcional ao número de usuários ativos.}~\cite{Kleppmann2017},
sobre a mesma pilha tecnológica Node.js/PostgreSQL. O \textbf{Cenário~A} é gerado integralmente
pelo Lovable a partir de prompts de negócio sem qualquer especificação arquitetural,
espelhando os padrões P1--P3.

O \textbf{Cenário~B} não é uma nova geração independente a partir dos mesmos requisitos
de negócio, mas uma refatoração do código funcionalmente equivalente ao Cenário~A
(Seção~\ref{sec:limitacoes}, item 5), aplicada via prompts arquiteturais explícitos
enviados a um assistente de codificação baseado em LLM (Claude, claude-sonnet-4-5),
seguindo o modelo de quatro componentes proposto na Seção~\ref{subsec:modelo_prompt}:
(i)~JOIN eliminando o padrão N+1; (ii)~pool otimizado (\texttt{max:20}, timeout 2\,s,
\textit{fail-fast}\footnote{Estratégia de falhar imediatamente ao detectar indisponibilidade, retornando 503 ao cliente em vez de deixá-lo suspenso.} 503); (iii)~\textit{cache} em memória com TTL
(\textit{time-to-live}) de 60--300\,s; e (iv)~fila assíncrona de escritas.

Essa escolha metodológica — refatorar em vez de gerar novamente do zero — foi
deliberada: manter o código-base do Cenário~A como ponto de partida do Cenário~B
garante que o escopo funcional (rotas, telas, comportamento) seja idêntico entre os
dois cenários, isolando a variável de interesse (padrão de acesso a dados) de
variações funcionais que uma nova geração independente poderia introduzir. Ao mesmo
tempo, essa decisão introduz duas variáveis não controladas entre os cenários: a
ferramenta de geração (Lovable, um construtor no-code, no Cenário~A; um assistente
de codificação baseado em LLM, no Cenário~B) e o modo de construção (geração autônoma
vs.\ refatoração guiada). Essa escolha reflete o próprio objeto de estudo: o modelo de
prompt arquitetural proposto (Seção~\ref{subsec:modelo_prompt}) foi desenhado para uso
com assistentes de codificação capazes de interpretar restrições arquiteturais
explícitas, não com construtores no-code como o Lovable, que abstraem essas decisões
do fundador (Seção~\ref{sec:rmodp}). A pilha tecnológica de destino
(Node.js/PostgreSQL), os dados e a infraestrutura de teste são idênticos entre os
cenários; a diferença de desempenho observada é atribuível majoritariamente às decisões
arquiteturais aplicadas, mas as variáveis de ferramenta e modo de construção não podem
ser completamente descartadas como fatores contribuintes, o que é discutido como
limitação na Seção~\ref{sec:limitacoes}.

\subsection{Protocolo Experimental}

Ambos os cenários foram avaliados pelo ATAM~\cite{Kazman2000}, com SLA (\textit{Service Level Agreement}) definido em latência $<$\,2.000\,ms (p95) e erro $<$\,1\%. Para isso, os testes de carga~\cite{grafana2024} foram executados com 100 usuários virtuais (VUs) simultâneos por 3 minutos na rota \texttt{GET /app/team}, em ambiente GitHub Codespaces (2\,vCPUs, 8\,GB RAM), com banco de dados PostgreSQL, servidor Node.js e K6 no mesmo host; cada usuário virtual realiza, em sequência, autenticação, acesso à trilha e registro de progresso.

Para descartar a hipótese de que a diferença de desempenho decorreria simplesmente de código mais complexo no Cenário~B, calculou-se também a Complexidade Ciclomática (CC)~\cite{McCabe1976} média por função em ambos os servidores, via a ferramenta \textit{lizard}. O resultado foi CC média de 2,4 no Cenário~A (15 funções, 231 linhas) contra 2,2 no Cenário~B (18 funções, 263 linhas) — ambas bem abaixo do limiar de risco (CC $>$ 15) e, portanto, estatisticamente equivalentes. Isso confirma que a diferença de desempenho não decorre de complexidade de código, mas do padrão de acesso a dados escolhido. Métricas de acoplamento e coesão orientadas a objetos (CBO, LCOM)~\cite{Chidamber1994} não se aplicam a este contexto, já que ambos os servidores seguem estilo procedural baseado em rotas Express, sem classes. Os dados completos dessa análise estão disponíveis no repositório de artefatos~\cite{GenAIDebt2026}.

%% ─────────────────────────────────────────────────────────────────────────────
\section{Cenário A: MVP Gerado por IAG}
%% ─────────────────────────────────────────────────────────────────────────────

O Cenário~A foi gerado pelo Lovable via prompt: ``Crie uma plataforma de onboarding
com banco de dados, login, trilhas de aprendizado, módulos com texto, barra de
progresso e painel do gestor'', sem parâmetros arquiteturais. Um achado
metodologicamente relevante: a resposta ``Admin + Gestor + Colaborador'' às perguntas
de clarificação da ferramenta gerou automaticamente RBAC (\textit{Role-Based Access
Control}) completo com RLS (\textit{Row-Level Security}) em 9 tabelas e três
\textit{security\- definer\- functions}, sem instrução explícita, evidenciando
elicitação implícita de requisitos funcionais pela IAG, com qualidade arquitetural
variável nos demais atributos de qualidade~\cite{ISO25010}.

\textbf{Pontos de sensibilidade identificados pelo ATAM.} Três problemas se reforçam mutuamente no Cenário~A. Primeiro, o dashboard do gestor executa múltiplas consultas ao banco em loop síncrono para cada colaborador, caracterizando o padrão N+1. Esse padrão, por sua vez, agrava um segundo problema: o pool de conexões, já subdimensionado e sem tempo limite configurado, satura rapidamente sob requisições paralelas, pois cada requisição N+1 mantém várias conexões abertas simultaneamente. Por fim, todas as escritas são síncronas e bloqueantes, competindo pelas mesmas conexões já disputadas pelas leituras. A análise ATAM sintetiza essa cadeia em três compromissos arquiteturais: N+1 vs.\ simplicidade do código; ausência de \textit{cache} vs.\ consistência de dados; e escritas síncronas vs.\ capacidade de atendimento sob pico.

A pilha tecnológica gerada autonomamente pelo Lovable compreendeu: TanStack Start como
framework de frontend, ambiente de execução (\textit{runtime}) Bun, implantação (\textit{deploy}) via Cloudflare Workers em funções
\textit{stateless} com restrições de memória, e Supabase (PostgreSQL) acessado
via SDK sem camada intermediária de servidor de aplicação. Arquiteturalmente, a pilha tecnológica não incluiu
mecanismos de cache nem processamento assíncrono: todas as operações são síncronas
e bloqueantes, decisões que refletem os padrões P1 e P3 da persona Onboard e
que se revelaram determinantes para o colapso em escala.

%% ─────────────────────────────────────────────────────────────────────────────
\section{Cenário B: Arquitetura Refatorada}
%% ─────────────────────────────────────────────────────────────────────────────

\subsection{Como Identificar Lacunas Arquiteturais sem Formação Técnica}

Fundadores identificam dívida arquitetural pelos sintomas observáveis em produção,
não pelo código-fonte. A Tabela~\ref{tab:sinais} apresenta os quatro sinais de
alerta mapeados na persona Onboard, cada um correspondendo a uma decisão
arquitetural corretiva específica no modelo de prompt proposto.

\begin{table}[t]
  \caption{Sinais de alerta de lacunas arquiteturais observáveis por fundadores sem
           formação técnica, mapeados aos componentes R1--R4.}
  \label{tab:sinais}
  \small
  \begin{tabular}{p{0.7cm}p{2.2cm}p{2.2cm}p{1.3cm}}
    \toprule
    \textbf{Sinal} & \textbf{O que o fundador observa} & \textbf{Indicação técnica} & \textbf{Ref.} \\
    \midrule
    S1 & Aplicação fica mais lenta à medida que dados crescem
       & Padrão N+1: mais registros $\Rightarrow$ mais queries & R1 \\
    S2 & Erros ``servidor indisponível'' com vários usuários simultâneos
       & Pool de conexões saturado & R2, R4 \\
    S3 & Dashboard do gestor lento mesmo com poucos colaboradores
       & Ausência de cache: tudo recalculado a cada requisição & R3 \\
    S4 & Marcar módulo como concluído demora vários segundos
       & Escrita síncrona bloqueia resposta HTTP (\textit{HyperText Transfer Protocol}) & R4 \\
    \bottomrule
  \end{tabular}
\end{table}

\subsection{O Modelo de Prompt Arquitetural para Fundadores}
\label{subsec:modelo_prompt}

O Cenário~B foi construído aplicando quatro refatorações ao servidor de aplicação do Cenário~A,
cada uma via prompt explícito ao Claude (claude-sonnet-4-5). Vogelsang~\cite{Vogelsang2024}
argumenta que requisitos precisam ser enriquecidos com decisões de design antes de
serem usados como prompts, trabalho de Engenharia de Requisitos irredutível à
automação. O modelo proposto operacionaliza essa decomposição em linguagem acessível
a fundadores sem formação técnica, como um \textit{architectural constraint prompt
pattern}~\cite{White2023}.

Na prática, o fundador reconhece uma funcionalidade lenta por um destes sinais: S1, lentidão progressiva conforme os dados crescem, indicando N+1; S2, \textit{timeouts} e erros sob pico de usuários, indicando pool saturado; ou S3, telas lentas mesmo com poucos dados, indicando ausência de \textit{cache}. É importante notar que o modelo é projetado exclusivamente para avaliar e corrigir escalabilidade em funcionalidades já existentes, não para orientar a construção de funcionalidades novas. A Tabela~\ref{tab:modelo_prompt} apresenta os quatro componentes do modelo, cujo template completo está disponível no repositório~\cite{GenAIDebt2026}. Vale destacar ainda que o componente~(4) cumpre uma função adicional de transferência de conhecimento: ao solicitar que a IAG explique o padrão corrigido, o fundador acumula compreensão técnica de forma iterativa a cada refatoração~\cite{Vogelsang2024}.

\begin{table}[t]
  \caption{Template de prompt arquitetural de quatro componentes para avaliação
           de escalabilidade de MVPs. O fundador substitui os campos entre colchetes
           pela descrição do seu próprio sistema.}
  \label{tab:modelo_prompt}
  \small
  \begin{tabular}{p{1.0cm}p{2.5cm}p{2.8cm}}
    \toprule
    \textbf{Comp.} & \textbf{O que escrever} & \textbf{Exemplo preenchido} \\
    \midrule
    (1) Função & Descreva a tela ou endpoint lento & ``Painel do gestor com lista de colaboradores e progresso'' \\
    (2) Restrição & Como o banco deve ser acessado (JOIN, cache, async) & ``Use uma única query SQL (\textit{Structured Query Language}) com JOINs, não loops'' \\
    (3) Tradeoff & Consequência aceita em troca de escalabilidade & ``Aceito dados com 60\,s de atraso em trilhas'' \\
    (4) Docs & Peça explicação do padrão corrigido e próximos passos & ``Comente o que mudaria para 10$\times$ mais usuários'' \\
    \bottomrule
  \end{tabular}
\end{table}

\subsection{As Quatro Refatorações e Seus Prompts}

Os prompts completos de R1–R4, com os quatro componentes preenchidos, estão disponíveis no repositório~\cite{GenAIDebt2026}. As quatro refatorações a seguir foram aplicadas em sequência, cada uma resolvendo um gargalo específico identificado na Seção~4 e habilitando a correção da próxima.

\subsubsection{R1: Eliminação do Padrão N+1 (Sinais S1, S3)}

A primeira refatoração ataca a causa raiz identificada na Seção~4: a rota \texttt{GET /app/team} executava $3+2M$ queries por requisição em um loop \texttt{for/await} síncrono, o que gerava 103 queries para $M=50$ colaboradores. \textbf{Decisão arquitetural:} o loop foi substituído por uma única query SQL com \texttt{LEFT JOIN}s e \texttt{GROUP BY}, retornando todos os colaboradores com seu progresso em uma só ida ao banco. \textbf{Resultado:} a contagem caiu de 103 para 1 query por requisição ($-99\%$), com comportamento $O(1)$ independente de $M$ — eliminando, na origem, a pressão sobre o pool de conexões que a próxima refatoração ainda precisa endereçar.

\subsubsection{R2: Pool de Conexões Otimizado (Sinal S2)}

Mesmo com o N+1 eliminado, o pool de conexões original (\texttt{max:10}, sem \textit{timeout}) ainda saturava com aproximadamente 10 requisições paralelas. \textbf{Decisão arquitetural:} o pool foi ampliado para \texttt{max:20}, com \textit{timeout} de 2\,s e retorno imediato de 503 (\textit{fail-fast}) quando esgotado, em vez de manter requisições penduradas indefinidamente. \textbf{Resultado:} a capacidade do pool dobrou (10$\rightarrow$20 conexões) e requisições em excesso passaram a falhar rapidamente e de forma sinalizada, em vez de ficarem suspensas sem resposta.

\subsubsection{R3: Cache em Memória com TTL (Sinal S3)}

Além do dashboard do gestor, endpoints quase-estáticos como \texttt{GET /app/tracks} e \texttt{GET /app/modules/:id} também iam ao banco a cada requisição, mesmo sem mudança de dados entre chamadas. \textbf{Decisão arquitetural:} foi adicionado \textit{cache} em memória (\texttt{node-cache}) com TTL de 60\,s para trilhas e 300\,s para módulos, aceitando consistência eventual nesse intervalo; rotas de progresso individual foram deliberadamente excluídas do \textit{cache}, por exigirem dados sempre atualizados. \textbf{Resultado:} taxa de acerto (\textit{cache hit}) esperada acima de 98\% para trilhas e 99\% para módulos, reduzindo a carga sobre o banco nesses endpoints a praticamente zero.

\subsubsection{R4: Fila Assíncrona de Escritas (Sinal S2)}

A última refatoração trata a competição por conexões entre leituras e escritas, apontada na Seção~4: o endpoint \texttt{POST /app/modules/:id/complete} executava um \textit{upsert}\footnote{Operação de banco de dados que insere um novo registro se não existir, ou atualiza o existente; no Cenário~A, executada de forma síncrona, bloqueando o HTTP até a gravação completar.} síncrono, bloqueando a resposta HTTP até a gravação ser concluída. \textbf{Decisão arquitetural:} a escrita passou a ser enfileirada em memória, com retorno imediato de \texttt{\{success:true, queued:true\}} — aceitando perda de registros em caso de reinicialização do servidor, em troca de latência de resposta próxima de zero. \textbf{Resultado:} o endpoint responde imediatamente e libera a conexão do pool, com consistência eventual de até 500\,ms, um custo aceitável para o domínio de onboarding corporativo.

Em conjunto, as quatro refatorações atuam em camadas complementares: R1 elimina o N+1 na origem; R2 converte esse ganho em capacidade real de atendimento; R3 elimina queries repetidas em rotas quase-estáticas; e R4 desacopla as escritas do pool de conexões. A seção seguinte avalia, via ATAM, os compromissos residuais que essa combinação de decisões introduz.

\subsection{Avaliação ATAM do Cenário B}

Dois pontos de sensibilidade residuais foram introduzidos: (1)~consistência eventual
de R3 e R4 (janelas de 60--500\,ms); e (2)~durabilidade da fila in-process de R4
(perda em reinicializações do servidor). Ambos foram avaliados como aceitáveis para
o domínio de onboarding corporativo: a latência de progresso de 500\,ms é
imperceptível ao usuário, e perda de registro em crash afeta menos de 0,1\% das
sessões sob carga estável. Os pontos residuais são documentados no código como
comentários de DTA, substituíveis por BullMQ + Redis quando o SLA de durabilidade
justificar o custo operacional. Os scripts K6, Dockerfiles e prompts completos estão
disponíveis no repositório público~\cite{GenAIDebt2026}, permitindo a
reprodução integral do experimento.

%% ─────────────────────────────────────────────────────────────────────────────
\section{Resultados}
%% ─────────────────────────────────────────────────────────────────────────────

Os resultados confirmam a hipótese: sistemas gerados por IAG sem especificação
arquitetural incorporam padrões que comprometem a escalabilidade. No Cenário~A,
a rota \texttt{GET /app/team} executa múltiplas consultas em loop por colaborador
(padrão N+1\footnote{Loop em JavaScript que gera uma consulta ao banco por colaborador, causando degradação proporcional ao número de registros.}); no Cenário~B, a mesma rota usa um único LEFT JOIN.
Ambos os cenários completaram 100 VUs por 3 minutos sem erros HTTP. Isso não
significa desempenho equivalente: o Cenário~A respondeu em média em 1.880\,ms,
ultrapassando o limite de desempenho definido (p95~$<$~2\,s), enquanto o
Cenário~B respondeu em 30\,ms, cumprindo todos os critérios. Sob cargas maiores, o
padrão N+1 do Cenário~A levaria ao esgotamento do pool de conexões, o mesmo
mecanismo documentado em produção pelos fundadores entrevistados. A diferença de
98,4\% no tempo de resposta é atribuível majoritariamente às decisões arquiteturais
aplicadas sobre pilha tecnológica, dados e infraestrutura idênticos; a ferramenta de
geração e o modo de construção do Cenário~B (refatoração guiada, e não geração
independente) constituem variáveis não controladas, discutidas na Seção~\ref{sec:limitacoes}.
A Tabela~\ref{tab:resultados} consolida os resultados; os artefatos estão no
repositório~\cite{GenAIDebt2026}.

\begin{table}[t]
  \caption{Desempenho em runtime dos Cenários~A e~B sob condições idênticas de teste
           (mesma pilha de destino Node.js/PostgreSQL, mesmos dados, mesma infraestrutura,
           100 VUs, 3 min). Ferramenta de geração e modo de construção do código-fonte
           diferem entre os cenários (ver linhas de contexto e Seção~\ref{sec:limitacoes}).}
  \label{tab:resultados}
  \small
  \begin{tabular}{lp{2.0cm}p{2.0cm}}
    \toprule
    \textbf{Contexto de construção} & \textbf{Cenário A} & \textbf{Cenário B} \\
    \midrule
    Ferramenta de geração    & Lovable (no-code) & Claude (refatoração guiada) \\
    Modo de construção       & Geração autônoma & Refatoração do código do Cenário~A \\
    \midrule
    \textbf{O que foi medido} & \textbf{Cenário A} & \textbf{Cenário B} \\
    \midrule
    Erros HTTP               & 0,00\%      & 0,00\%          \\
    Tempo de resposta médio  & 1.880\,ms   & 30\,ms          \\
    Limite p95 $<$ 2\,s      & \textbf{violado} & cumprido   \\
    \midrule
    Consultas por requisição & 109 (em loop) & 1 (JOIN)      \\
    Escrita de dados         & Síncrona    & Assíncrona      \\
    \textit{Cache}           & Ausente     & Ativo (TTL 60\,s) \\
    \bottomrule
  \end{tabular}
\end{table}

\subsection{Discussão}

A análise estática confirma a hipótese central: prompts sem especificação
arquitetural geram sistemas funcionalmente corretos, mas com decisões de
implementação que comprometem a escalabilidade~\cite{Shypula2023,Bass2021}. A
diferença entre os cenários é verificável diretamente no código: 109 consultas
ao banco por requisição no Cenário~A contra 1 no Cenário~B. Essa diferença, de
$O(N)$ para $O(1)$, é a causa raiz da degradação observada nos testes de carga,
e é inspecionável por qualquer desenvolvedor que examine o repositório.

O achado mais relevante é a \textbf{desproporção entre esforço e impacto}:
94 linhas de código a mais (25\%) e cerca de uma hora de prompts arquiteturais
produziram redução de 98,4\% na latência média. Investir em especificação
arquitetural desde o início é muito mais barato do que refatorar um sistema
que já apresenta problemas em produção.

Um detalhe importante do modelo proposto é o papel da componente~(3), o compromisso (\textit{tradeoff})
explícito. Ao escrever frases como \textit{``aceito 503 imediato''} ou
\textit{``aceito dados com 60\,s de atraso''}, o fundador dá ao modelo de IAG
a permissão para escolher a solução que escala, em vez da que parece mais simples.
Sem essa declaração, modelos de IAG tendem a gerar código que prioriza consistência
imediata~\cite{Nascimento2024}, o que é razoável em tutoriais, mas problemático
sob carga real.

Para a comunidade de Engenharia de Software, o estudo aponta para uma mudança de
perspectiva: o problema com sistemas gerados por IAG não é de correção sintática,
mas de especificação de requisitos não-funcionais. O que falta não está nas
ferramentas, está no prompt.

Os resultados sugerem que a principal limitação da geração automática de código por
IAG não reside na capacidade de produzir software
funcional, mas na ausência de mecanismos capazes de incorporar requisitos
arquiteturais explícitos durante o processo de geração.

%% ─────────────────────────────────────────────────────────────────────────────
\section{Considerações Finais}
%% ─────────────────────────────────────────────────────────────────────────────

\subsection{Contribuições e Implicações Práticas}

A principal contribuição empírica é a quantificação da dívida técnica arquitetural introduzida por IAG:
o Cenário~A executa 109 consultas ao banco por requisição em sua rota crítica, contra
1 no Cenário~B, diferença inspecionável no código-fonte e corroborada pelos testes de
carga, que evidenciaram latência média de 1.880\,ms vs.\ 30\,ms (−98,4\%) com 100
VUs simultâneos, e vazão (\textit{throughput}) de 29 vs.\ 102\,req/s (+250\%).
A contribuição científica original é o \textbf{modelo de prompt arquitetural de
quatro componentes}, que operacionaliza a especificação de requisitos não-funcionais
de escalabilidade em linguagem acessível a fundadores sem formação
técnica~\cite{White2023,Nascimento2024,Vogelsang2024}. O modelo preenche uma lacuna
na interseção de \textit{prompt engineering} estruturado, elicitação de requisitos
não-funcionais e arquitetura de software assistida por IAG. Para a indústria, a
lacuna identificada é de especificação de requisitos, não de capacidade das
ferramentas.

\subsection{Limitações}
\label{sec:limitacoes}

Seis limitações principais: (1)~os testes de carga foram executados em GitHub
Codespaces com banco de dados, servidor e K6 no mesmo host, o que impede isolar
o gargalo arquitetural do gargalo de infraestrutura sob cargas elevadas; os
resultados dinâmicos devem ser interpretados como evidência qualitativa da
degradação, não como métricas absolutas; (2)~dados de negócio sintéticos, que
podem subrepresentar o impacto do N+1 com volumes reais; (3)~fundamento em uma
única startup, restringindo a generalização dos padrões P1--P5 a outros domínios
e tecnologias; (4)~avaliação ATAM e RM-ODP realizada por um único autor, de
caráter qualitativo-comparativo; e (5)~por limitações de instrumentação de carga
sobre a pilha serverless original (Cloudflare Workers/Supabase), o Cenário~A
submetido ao teste de carga é uma reprodução funcionalmente equivalente do
padrão N+1 observado no aplicativo gerado pelo Lovable, implementada em
Node.js/PostgreSQL sobre a mesma pilha do Cenário~B para permitir comparação
controlada. O aplicativo originalmente gerado pelo Lovable (TanStack
Start/Bun/Cloudflare Workers/Supabase) está disponível no repositório de
artefatos para inspeção do padrão arquitetural original~\cite{GenAIDebt2026}; e
(6)~o Cenário~B não foi gerado de forma independente a partir dos mesmos requisitos
de negócio do Cenário~A, mas obtido por refatoração do código deste, aplicada por
uma ferramenta de geração distinta (assistente de codificação baseado em LLM, em vez
do construtor no-code usado no Cenário~A). Ferramenta de geração e modo de construção
(geração autônoma vs.\ refatoração guiada) são, portanto, variáveis não controladas
entre os cenários, o que impede atribuir a diferença de desempenho exclusivamente ao
nível de especificação arquitetural do prompt -- a interpretação mais conservadora é
que os resultados demonstram o impacto de decisões arquiteturais explícitas aplicadas
a um sistema já funcional, não necessariamente o impacto isolado do nível de
especificação em uma geração original. Trabalhos futuros devem replicar o Cenário~B
como uma geração independente pela mesma ferramenta do Cenário~A, usando prompts com
e sem especificação arquitetural, para isolar essa variável.

\subsection{Trabalhos Futuros}

Quatro direções: (1)~replicação do experimento em ambiente com instâncias
separadas (servidor, banco e ferramenta de carga em hosts distintos), que permitiria
isolar o gargalo arquitetural e obter métricas de VUs e latência com maior validade
externa; (2)~extensão do protocolo de entrevistas para identificar P6--P10 em
startups de diferentes setores e tecnologias; (3)~desenvolvimento de um catálogo de
prompts arquiteturais cobrindo autenticação, observabilidade e resiliência; e
(4)~investigação de abordagens automatizadas para detecção de DTA em repositórios
gerados por IAG via análise estática combinada com testes de carga em CI/CD (\textit{Continuous Integration/Continuous Delivery}).

Este trabalho indica que o principal desafio da Engenharia de Software na era da
IAG não consiste apenas em desenvolver modelos mais
capazes de gerar código, mas em tornar explícitos os requisitos arquiteturais que
orientam esse processo de geração, reduzindo a introdução inadvertida de dívida
técnica arquitetural desde as primeiras versões do software.

%% ─────────────────────────────────────────────────────────────────────────────
\section*{Disponibilidade de artefatos}
%% ─────────────────────────────────────────────────────────────────────────────

Os artefatos deste estudo estão disponíveis em repositório
público~\cite{GenAIDebt2026},
incluindo: código-fonte completo dos Cenários~A e~B (\texttt{server.js},
\texttt{schema.sql}, \texttt{seed.js}, Dockerfiles), scripts K6 para reprodução
dos testes de carga, os quatro prompts arquiteturais completos de R1--R4 em
\texttt{prompts/\linebreak[0]prompts-arquiteturais.md}, a análise de complexidade
ciclomática em \texttt{analise-estatica/\linebreak[0]complexidade-ciclomatica.md},
e o código-fonte original gerado pelo Lovable (pasta \texttt{src/}, stack
TanStack~Start/Supabase/Cloudflare~Workers). O README contém instruções de
reprodução e observações sobre a configuração de ambiente recomendada (instâncias
separadas para banco, servidor e K6). Os dados qualitativos da entrevista não são
disponibilizados em cumprimento ao compromisso de anonimização (TCLE).

%% ─────────────────────────────────────────────────────────────────────────────
\section*{Agradecimentos}
%% ─────────────────────────────────────────────────────────────────────────────

A primeira autora agradece ao Instituto de Tecnologia e Liderança (INTELI) pela
bolsa de extensão de Iniciação Científica e pela oportunidade de desenvolver esta
pesquisa. Agradece ao Prof.\ Reginaldo Arakaki e à Prof.\textsuperscript{a} Fabiana
Martins de Oliveira pela orientação dedicada ao longo de toda a pesquisa: as reuniões
constantes, os feedbacks detalhados a cada versão e a rotina de leitura crítica
compartilhada foram fundamentais para a qualidade deste trabalho. Agradece também
aos fundadores da startup pela disponibilidade, pela generosidade em compartilhar
sua experiência e pelo consentimento com o uso dos dados para fins acadêmicos.

Em conformidade com a política de uso de IAG do VII~CTIC-ES, declara-se
que o Claude (Anthropic) foi utilizado para auxílio na geração das tabelas em
\LaTeX{} e do gráfico de radar (Figura~\ref{fig:rmodp}). Todo o conteúdo científico,
incluindo hipóteses, metodologia, implementação dos cenários, execução dos testes e
interpretação dos resultados, é de autoria exclusiva das pesquisadoras.

%% ─────────────────────────────────────────────────────────────────────────────
\bibliographystyle{ACM-Reference-Format}
\bibliography{referencias}
%% ─────────────────────────────────────────────────────────────────────────────

\end{document}
